\subsection{C.28 at V ms.~46--49: stages, unions, and strict retractions}
Call a category $A$ local when the specified comparison
$\pi_0\operatorname{Fun}(X,A)\to[BX,BA]$ is bijective for every small
$X$. This is stronger than Thomason fibrancy. The model structure used
here is the one stated and proved in
\cite[Theorem~2.3]{may2017posets}: its generating trivial cofibrations
are $c\operatorname{Sd}^2\Lambda_i^r\to c\operatorname{Sd}^2\Delta^r$,
and fibrancy means that $\operatorname{Ex}^2NA$ is Kan. These generating
maps have finite source and target and are full sieve inclusions: the
twice-subdivided simplex is the poset of chains of faces, and a
subcomplex gives the downward-closed subposet of its face chains.
These are exactly the hypotheses used in the following checks.

\begin{proposition}[The positive presentation assertions in C.28]
\label{F:C28-presentations}
A globally invariant sequential ind-object admits a presentation by
full inclusions of local Thomason-fibrant categories. The union is
Thomason-fibrant. It need not be local or represent the ind-object.
Filtered colimits of Thomason-fibrant categories are fibrant, and
objectwise weak equivalences of filtered diagrams induce weak
equivalences of their colimits.
\end{proposition}
\begin{proof}
\textbf{The groupoid presentation.} Theorem~\ref{F:global-ind-groupoids}
gives the specified ind-isomorphism to the actual groupoid sequence
$\Gamma_n=\Pi_1K_n$. Let $D_n$ be its covariant Grothendieck construction
over $\{0<\cdots<n\}$. Its objects are $(i,x)$ with $x\in\Gamma_i$;
for $i\leq j$ its arrows $(i,x)\to(j,y)$ are the arrows
$\Gamma_{ij}(x)\to y$ of $\Gamma_j$. The product of labels is
$h\,\Gamma_{jk}(g)$, which verifies associativity and identities.
The terminal-fibre retraction $r_n:D_n\to\Gamma_n$ sends objects and
labels forward. Its section $j_n$ satisfies $r_nj_n=1$, and identity
labels give $1\Rightarrow j_nr_n$. Thus $D_n$ is categorically homotopy
equivalent to a groupoid and is local: functors to a groupoid factor
through fundamental groupoids, and their transformations are precisely
the natural isomorphisms classifying unpointed maps of classifying
spaces. Applying the two inverse zigzags transports this comparison to
$D_n$.

\textbf{Fullness and fibrancy.} The $D_n$ are full subcategories of
$D_{n+1}$, with strictly compatible retractions. Given a generating
horn $A\subset B$ and $v:A\to D_n$, first extend $r_nv$ to
$w:B\to\Gamma_n$. A groupoid has Kan nerve and hence is
Thomason-fibrant, so this extension exists. Retain $v$ on old objects
and send each new object to its value under $w$ in the terminal fibre.
Retain old arrows and label all new arrows by $w$. Since $A$ is a full
sieve, there is no arrow from a new object to an old one; all other
new arrows have allowed terminal-fibre targets. Mixed compositions
are exactly multiplication of their labels in $\Gamma_n$, so the
extension is a strict functor. This checks the lifting assertion in
the written proof of C.28.

\textbf{Filtered passages.} A lifting problem for a generating horn
in a filtered colimit is finite data and finite equations, so it occurs
at a stage and is solved there. This proves fibrancy of the colimit.
Nerves commute with filtered colimits by the same finite-data argument;
their ordinary colimits agree weakly with homotopy colimits as checked
in Proposition~\ref{F:C22-types}. Objectwise weak equivalences therefore
induce weak equivalences of these colimits. For an arbitrary actual
sequence, compatible fibrant replacements also exist: choose trivial
cofibrations $K_n\to R_n$ to fibrant targets, and extend
$K_n\to K_{n+1}\to R_{n+1}$ across $R_n$ by lifting. A sequence is
freely generated by its adjacent arrows, so no coherence equations
remain to be solved. None of these lifting arguments asserts locality
of the union. The next counterexample verifies its failure.
\end{proof}

The previously missing all-type path-space comparison is now proved
in Theorem~\ref{F:unmarked-path-comparison}. Its relative universal
diagrams supply the coherence in every homotopy degree and check the
specified sequential span. The independent C.20 component proof,
C.26 diagram test and C.28 finite-test theorem do not need that repair.

\begin{proposition}[The union and uniqueness counterexample in C.24/C.28]
\label{F:C28-union-counterexample}
There is a sequence of finite local fibrant categories with full
inclusions whose fibrant union $T$ is not local. It has an endofunctor
$s$ equal to the identity as an ind-morphism in
$\operatorname{Ind}(\Cat_{\mathrm{hot}})$ but not a nerve weak
equivalence. Thus ind-isomorphisms do not give the asserted uniqueness
of unions at V ms.~46--49.
\end{proposition}
\begin{proof}
\textbf{The tables and stage retractions.} Let
$G_n=\operatorname{Sym}\{0,\ldots,n\}$ and take $T_n$ with objects
$0,\ldots,n$, with $\Hom(i,j)=G_j$ for $i\leq j$, and with no arrows
otherwise. Compose labels by
$(h:j\to k)(g:i\to j)=h\,\iota_{jk}(g)$.
Group multiplication verifies every category equation. This is the
preceding $D_n$ for inclusions of one-object groupoids, so its specified
terminal-fibre section, retraction and identity-label transformation
make it local and fibrant; the transitions are full inclusions.

\textbf{The actual ind-homotopies.} Set $s(i)=i+1$ and shift the support
of every permutation one place, fixing zero. This commutes with label
composition. For fixed $n$, both the inclusion and the shift
$T_n\to T_{n+1}$ have identity-label transformations to their
corresponding functors constant on objects at $n+1$. The latter two
are naturally isomorphic by the cycle $c=(0\ 1\ \cdots\ n+1)$,
because $cgc^{-1}=s(g)$ for every $g\in G_n$.
The resulting three-arrow zigzag proves equality with the identity
ind-morphism. This is an actual zigzag at the next stage, not just an
equality of maps of classifying spaces.

\textbf{Failure on the union.} The stage retractions assemble to
$r:T\to\mathbf BG$, where
$G=\operatorname{Sym}_{\mathrm{fin}}(\N)$. The preceding filtered
comparison proves $Br$ a weak equivalence. Under this identification,
$Bs$ induces the group shift. Its image fixes zero and omits $(0\ 1)$,
so it is not a fundamental-group isomorphism and cannot be homotopic
to the identity. A functor $\mathbf BG\to T$ selects one object $i$
and has all its arrow images in the finite group $G_i$. Its classifying
map cannot invert $Br$, even up to unpointed homotopy, since conjugation
preserves finiteness of the image. Thus the local comparison for this
source fails. The weak equivalence $r$ between fibrant categories cannot
have a $\Cat_{\mathrm{hot}}$ inverse either.

\textbf{Finite verification and the other printed obstructions.}
The independent program \texttt{f\_union\_models.py} supplies all
tables through $T_3$, with arrow counts $1,5,23,119$. It checks every
identity and composable triple (respectively $1,23,917,73037$), each
retraction and all three transformations into the next stage.
For the separate unbounded-length obstruction, finite intervals of the
infinite alternating fence contract by deleting endpoints. Its full
union admits no finite categorical homotopy from identity to a constant:
a zigzag of length $l$ would move every vertex by undirected distance
at most $l$, while distances to a fixed vertex are unbounded. Finally
the terminal fork $0\to1\leftarrow2$ is fibrant: a functor on a full
sieve extends by the terminal object on new objects. Its vertices $0,2$
are connected but have no arrow in either direction. This separately
refutes the one-transformation inference.
\end{proof}

\begin{proposition}[The representation and retraction qualifications in C.28]
\label{F:C28-retractions}
For full-inclusion sequences, an ind-morphism in $\Cat$ is exactly a
functor on unions whose restriction to each source stage lands in some
target stage. Equality in $\operatorname{Ind}(\Cat_{\mathrm{hot}})$
does not imply equality or homotopy of the induced union functors.
A globally invariant ind-object is represented in $\Cat_{\mathrm{hot}}$
exactly when it is isomorphic to a constant ind-object.

If $A$ is Thomason-fibrant, $i:A\hookrightarrow B$ is a Thomason
cofibration, and $p:B\to A$ has $pi$ categorically homotopic to $1_A$,
there is a strict retraction $r$ in the same ordinary homotopy class as
$p$. If $A$ is also local, the equality holds in
$\Cat_{\mathrm{hot}}$. For a groupoid target the latter assertion
holds for any full inclusion injective on objects.
\end{proposition}
\begin{proof}
\textbf{Representation and bounded union maps.} The formula for these
morphisms is
\[
 \Hom_{\operatorname{Ind}(\Cat)}((A_n),(B_m))
    =\lim_n\colim_m\Hom(A_n,B_m).
\]
With literal full inclusions, equality
in the colimit means equality of the corresponding union maps. Compatible
representatives therefore glue to one functor on the union; conversely
a functor satisfying the displayed boundedness yields those compatible
classes. Replacing Hom by categorical homotopy classes only provides
eventual zigzags. Proposition~\ref{F:C28-union-counterexample} verifies
their failure to glue even to an ordinary homotopy of union maps.
The constant ind-embedding is fully faithful, and its associated
presheaf is the representable one. Yoneda gives the claimed exact
representation criterion. A weak equivalence between local categories
has a categorical-homotopy inverse: take its inverse in $\mathrm{Hot}$,
represent it using locality of its target, and use injectivity of both
local comparisons for the two composite equations. The preceding $r$
is an explicit failure when only fibrancy is assumed.

\textbf{The path-object construction, checking INDLOCAL D.1.}
Represent $pi\simeq1_A$ by a functor $A\times Z\to A$ for a finite
oriented interval $Z$ with endpoints $0,z$. The constant-path functor
$c:A\to\operatorname{Fun}(Z,A)$ is a categorical homotopy equivalence:
successive endpoint deletions contract $Z$, and their transformations
induce the inverse zigzags on functor categories. Factor endpoint
evaluation as a trivial cofibration followed by a fibration,
\[
 \operatorname{Fun}(Z,A)\xrightarrow{j}P
       \xrightarrow{(d_0,d_1)}A\times A.
\]
Since $A$ is fibrant, both projections of $A\times A$ are fibrations,
so $d_0,d_1$ are fibrations. The map $jc$ is a weak equivalence and
$d_0jc=d_1jc=1$, so both are trivial fibrations. The adjoint of the
given homotopy followed by $j$ gives $h:A\to P$, with $d_0h=pi$ and
$d_1h=1$. Lifting against $d_0$ extends it to
$\widetilde h:B\to P$ with $d_0\widetilde h=p$. Put
$r=d_1\widetilde h$. Then $ri=1$ strictly. In $\mathrm{Hot}$ both
$d_0,d_1$ equal the inverse of $jc$, so $[r]=[p]$ there. No finite
natural-transformation zigzag on $B$ follows from this factorization
alone.

\textbf{The two upgrades, checking INDLOCAL D.2.} Locality makes
$\pi_0\operatorname{Fun}(B,A)\to[BB,BA]$ injective and supplies the
first upgrade. For an actual groupoid, compose the given zigzag into
an isomorphism $\alpha:pi\Rightarrow1_A$. Identify $A$ with its full
image. Put $r(b)=b$ on $A$ and $r(b)=p(b)$ outside; use
$\beta_b=\alpha_b$ on $A$ and identity outside. On an arrow $u:b\to b'$
define $r(u)=\beta_{b'}p(u)\beta_b^{-1}$. Cancellation checks composition,
naturality of $\alpha$ gives $ri=1$, and $\beta:p\Rightarrow r$ is the
specified natural isomorphism. These check every qualification of the
retraction sentences used in C.28.
\end{proof}
